% !TeX root = ../main.tex
% 联动收益比较须使用相同预算、信息与需求条件。
% 两行动交互效应与联动策略相对分离策略的收益分别评价。
\chapter{道路抢通与资源调拨的联动决策机制}
\label{chap:coordination}

\section{道路恢复与资源交付的联动问题}
\subsection{抢通行动对通行能力的影响}
\subsection{资源调拨对服务恢复的作用}
\subsection{共同预算与行动时序约束}

\ThesisBlank

\section{抢通与调拨的统一动态决策模型}
\subsection{道路能力、库存与在途资源的状态转移}
\subsection{需求时限与服务损失评价}
\subsection{行动成本与联合可行域}

\ThesisBlank

\section{互补、替代与可分离条件分析}
\subsection{串联瓶颈下的抢通与调拨互补条件}
\subsection{替代通道与需求饱和下的边际收益变化}
\subsection{目标和约束可分解的充分条件}

\ThesisBlank

\section{联动策略与分离决策基准}
\subsection{抢通与调拨的联合规划及滚动更新}
\subsection{固定行动计划与分离规划基准}
\subsection{相同预算和信息条件下的可行性比较}

\ThesisBlank

\section{联动机制与收益来源的受控实验}
\subsection{抢通与调拨两因素交互检验}
\subsection{同预算下联动与分离策略的净收益比较}
\subsection{瓶颈位置与恢复时序的机制对照}

\ThesisBlank

\section{泸定情景验证与反例分析}
\subsection{不同灾损与资源条件下的策略表现}
\subsection{最优联动的零增益条件与近似策略的负收益来源}
\subsection{信息误差与抢通工期偏差的影响边界}

\ThesisBlank

\section{本章小结}
\ThesisBlank
